
Aufzählung sind auch sehr einfach:

\begin{enumerate}
    \item erster Punkt
    \item zweiter Punkt
    \begin{enumerate}
        \item erster Unterpunkt
        \item zweiter Unterpunkt
    \end{enumerate}
    \item weiter Punkt
\end{enumerate}

\blindtext



Kapitel beginnen mit \texttt{\\chapter}.

\section{Hintergrund} % (fold)
\label{sec:hintergrund}

\begin{center}
    ``Hübsches Zitat mit Referenz in der Fußnote.''\footnote{Autor Buchtitel Verlag, Erscheinungsjahr, Seite.}
\end{center}

Weitere Zitat \cite{weise2014benchmarking}. Und ein Conference Proceeding \cite{Hess2017}. Und ein Techreport \cite{instance1290}.

\section{Weitere Section} % (fold)
\label{sec:weitere_arbeiten}

\subsection{Untersection}
\label{ssec:untersection}

\blindtext


Tabellen, wie Tabelle \ref{tab:example} können auch mit einem Tabllengenerator wie auf Seite \url{https://www.tablesgenerator.com} im Internet erstellt werden.
\begin{table}[h]
    \centering
    \begin{tabular}{|r||l|l|}\hline
        t & $P(C_1)$ & $P(C_2)$\\\hline
        a & $0,6$ & $0,4$\\\hline
        b & $0,1$ & $0,9$\\\hline
        c & $0,82$ & $0,18$\\\hline\hline
        Mittelwert & \textbf{0,507} & $0,493$\\\hline
    \end{tabular}
    \caption{Beispiel für eine Tablle. Neben $Formeln$ und \textbf{Fettdruck} ist auch \textit{Kursiv} möglich.}
    \label{tab:example}
\end{table} 


Man kann auch \textbf{\textit{Fett und Kursiv}} kombinieren.

\begin{itemize}
    \item erster Punkt
    \item zweiter Punkt
    \begin{itemize}
        \item ein Unterpunkt
        \item noch ein Unterpunkt
    \end{itemize}
    \item ein weiterer Punkt
\end{itemize}